\chapter{Calcul et optimisation}\label{ch:iv1:calculus-and-optimisation}

\begin{ivtrains}
Calculer avec un esprit clair~: intégrales par symétrie, multiplicateurs de Lagrange, convexité, pas de gradient et de Newton, formule de Stirling, et les endroits où l'arithmétique en virgule flottante donne de mauvais résultats. Les questions attendent un raisonnement court et une vérification que le résultat est du bon ordre de grandeur.
\end{ivtrains}

\begin{question}[$\star$ \iqroles{QR, QT, Trader, Analyste} \iqfirm{toutes} \iqlevel{junior}]\label{iq:iv1:calculus-and-optimisation:1}
Où $x^x$ atteint-il son minimum pour $x > 0$, et quelle est la valeur de ce minimum~? Sans calculatrice, lequel des deux est le plus grand~: $e^\pi$ ou $\pi^e$~? Démontrez-le à l'aide d'une seule fonction.
\end{question}

\begin{question}[$\star$ \iqroles{QR, QT, DS, QD} \iqfirm{toutes} \iqlevel{junior}]\label{iq:iv1:calculus-and-optimisation:2}
Calculez $\int_{-\infty}^{\infty} e^{-x^2}\,dx$ en considérant son carré et en passant en coordonnées polaires. Utilisez ce résultat pour trouver la constante de la densité normale, puis l'espérance de $|Z|$ lorsque $Z$ suit une loi normale centrée réduite.
\end{question}

\begin{question}[$\star$ \iqroles{QR, QT, QD, Dev} \iqfirm{toutes} \iqlevel{junior}]\label{iq:iv1:calculus-and-optimisation:3}
Appliquez la méthode de Newton à $f(x) = x^2 - 2$, en partant de $x_0 = 1$. Écrivez la récurrence sous la forme d'une formule simple, puis donnez les trois premiers itérés sous forme de fractions. Combien de chiffres corrects chaque étape ajoute-t-elle~?
\end{question}

\begin{question}[$\star$ \iqroles{QR, QT, Trader, Analyste} \iqfirm{toutes} \iqlevel{junior}]\label{iq:iv1:calculus-and-optimisation:4}
Estimez $\ln 1.1$, $e^{0.1}$ et $\sqrt{1.04}$ à trois décimales près, en n'utilisant que deux ou trois termes d'une série de Taylor. Une obligation zéro-coupon, qui verse 1 au bout de 10 ans, est évaluée à $e^{-10y}$, où $y$ est son rendement. En utilisant une approximation au premier ordre puis au second ordre, quelle est la variation en pourcentage du prix lorsque le rendement augmente d'un point de pourcentage~?
\end{question}

\begin{question}[$\star\star$ \iqroles{QR, QT, Trader} \iqfirm{toutes} \iqlevel{junior}]\label{iq:iv1:calculus-and-optimisation:5}
Calculez $\int_0^{\pi/2} \frac{\sin^n x}{\sin^n x + \cos^n x}\,dx$ pour tout $n$ strictement positif. Calculez ensuite $\int_0^1 \frac{\ln(1 + x)}{1 + x^2}\,dx$, en utilisant une astuce analogue après le changement de variable $x = \tan\theta$.
\end{question}

\begin{question}[$\star\star$ \iqroles{QR, QT, DS, QD} \iqfirm{toutes} \iqlevel{junior}]\label{iq:iv1:calculus-and-optimisation:6}
Établissez la formule de Stirling pour $n!$ à partir de $n! = \int_0^\infty x^n e^{-x}\,dx$, en effectuant un développement autour du maximum de l'intégrande. Utilisez-la pour en donner une approximation de $\binom{100}{50}$. Quelle est la probabilité d'obtenir exactement 50 piles en 100 lancers~?
\end{question}

\begin{question}[$\star\star$ \iqroles{QR, QT, QD} \iqfirm{toutes} \iqlevel{junior}]\label{iq:iv1:calculus-and-optimisation:7}
Soit $I_n = \int_0^{\pi/2} \sin^n x\,dx$. Établissez une relation de récurrence par intégration par parties, puis calculez $I_{10}$ exactement. Pour $n$ grand, dans quelle mesure l'approximation $\sqrt{\pi/(2n)}$ est-elle précise~?
\end{question}

\begin{question}[$\star\star$ \iqroles{QR, QT, Analyste, DS} \iqfirm{toutes} \iqlevel{junior}]\label{iq:iv1:calculus-and-optimisation:8}
Maximisez $xyz$ sous la contrainte $x + y + z = 1$, à variables positives, à l'aide des multiplicateurs de Lagrange, puis par une autre méthode. Maximisez ensuite $\sqrt{x} + \sqrt{y}$ sous la contrainte $x + y = B$. Montrez que, pour $B = 2$, le multiplicateur est égal à la variation de la valeur optimale par unité de $B$, et expliquez ce que cela signifie.
\end{question}

\begin{question}[$\star\star$ \iqroles{QR, QT, DS, Analyste} \iqfirm{toutes} \iqlevel{junior}]\label{iq:iv1:calculus-and-optimisation:9}
Lesquelles de ces fonctions sont convexes~: $-\ln x$, $x^2$, $|x|$, $\sqrt{x}$, $\max(x, y)$, $xy$~? Donnez une raison en une ligne pour chacune. Utilisez l'inégalité de Jensen pour expliquer pourquoi la croissance composée d'un rendement annuel suivant une loi normale de moyenne 10~\% et de volatilité 20~\% se fait à un taux inférieur à 10~\%, puis estimez ce taux.
\end{question}

\begin{question}[$\star\star$ \iqroles{QR, DS, QD} \iqfirm{toutes} \iqlevel{junior}]\label{iq:iv1:calculus-and-optimisation:10}
La descente de gradient minimise $\frac12 x^\top Ax$, où $A$ admet pour valeurs propres 1 et 100. Quelles valeurs du pas assurent la convergence, et quelle est la meilleure~? Combien d'itérations faut-il pour réduire l'erreur d'un facteur $10^6$~? Que fait la méthode de Newton sur cette fonction, et que fait le momentum de type boule lourde~?
\end{question}

\begin{question}[$\star\star$ \iqroles{QR, QD, Dev, QT} \iqfirm{toutes} \iqlevel{junior}]\label{iq:iv1:calculus-and-optimisation:11}
Donnez un exemple dans lequel la méthode de Newton s'éloigne de la racine à chaque pas, quel que soit le point de départ. Pour $f(x) = \arctan x$, pour quelles valeurs initiales échoue-t-elle~? Pourquoi la résolution de la volatilité implicite demande-t-elle de la prudence, et comment sécuriseriez-vous une itération de Newton~?
\end{question}

\begin{question}[$\star\star$ \iqroles{QR, QD, Dev, LLE} \iqfirm{toutes} \iqlevel{junior}]\label{iq:iv1:calculus-and-optimisation:12}
En double précision, la plus petite racine de $x^2 - 10^8x + 1 = 0$ obtenue par la formule usuelle est fausse. Que donne-t-elle, quelle est la vraie valeur, et comment corrigeriez-vous ce calcul~? Pourquoi $1 - \cos x$ donne-t-il un mauvais résultat pour les petites valeurs de $x$, et par quoi le remplaceriez-vous~? Pourquoi la formule classique de la variance échoue-t-elle lorsque les données sont affectées d'un grand décalage, et quelle méthode préféreriez-vous~?
\end{question}

\begin{question}[$\star\star$ \iqroles{QR, QD, Dev} \iqfirm{toutes} \iqlevel{junior}]\label{iq:iv1:calculus-and-optimisation:13}
Vous estimez une dérivée par une différence avant en double précision. Quel est le pas optimal, et quelle précision atteignez-vous~? Comment ces deux grandeurs changent-elles pour une différence centrée~? Existe-t-il un moyen d'atteindre une précision presque complète pour une fonction analytique~?
\end{question}

\begin{question}[$\star\star\star$ \iqroles{QR, QT} \iqfirm{hedge fund, teneur de marché} \iqlevel{junior, confirmé}]\label{iq:iv1:calculus-and-optimisation:14}
Calculez $\int_0^{\pi/2} \ln \sin x\,dx$ sans recourir à aucune série, en utilisant seulement les symétries du sinus et la formule de l'angle double.
\end{question}

\begin{question}[$\star\star\star$ \iqroles{QR, QT} \iqfirm{hedge fund} \iqlevel{junior, confirmé}]\label{iq:iv1:calculus-and-optimisation:15}
Calculez $\int_0^\infty \frac{\sin x}{x}\,dx$ en introduisant un paramètre~: posez $I(t) = \int_0^\infty e^{-tx}\frac{\sin x}{x}\,dx$, dérivez sous l'intégrale, puis résolvez. Quelle est la valeur de $I(1)$~? Pourquoi l'intégrale d'origine converge-t-elle, sans toutefois converger absolument~?
\end{question}

\begin{question}[$\star\star\star$ \iqroles{QR, QT, DS} \iqfirm{hedge fund} \iqlevel{junior, confirmé}]\label{iq:iv1:calculus-and-optimisation:16}
Soit $A$ une matrice $n$ par $n$, symétrique définie positive. Calculez $\int_{\mathbb{R}^n} e^{-\frac12 x^\top Ax}\,dx$. Déduisez-en la covariance de la loi normale correspondante. Évaluez l'un et l'autre pour la matrice $A$ dont les lignes sont $(2, 1)$ et $(1, 2)$. Quel est le signe de la corrélation, et pourquoi~?
\end{question}

\begin{question}[$\star\star\star$ \iqroles{QR, QT, DS, Risque} \iqfirm{hedge fund} \iqlevel{junior, confirmé}]\label{iq:iv1:calculus-and-optimisation:17}
Vous allouez une unité de capital entre trois stratégies, avec des montants $x_i \ge 0$, afin de maximiser $\sum_i \ln(1 + x_i/n_i)$, où $n = (0.2, 0.5, 1)$. Écrivez les conditions d'optimalité, avec les contraintes de signe, puis résolvez. Pourquoi la dernière stratégie n'est-elle pas financée~? Comment la réponse change-t-elle si le capital est de 3 unités~?
\end{question}

\begin{question}[$\star\star\star$ \iqroles{QR, QD, DS} \iqfirm{toutes} \iqlevel{junior, confirmé}]\label{iq:iv1:calculus-and-optimisation:18}
Démontrez que la méthode de Newton converge quadratiquement au voisinage d'une racine simple. Montrez que la convergence n'est que linéaire pour une racine double, en prenant $f(x) = (x - 1)^2$, puis déterminez le taux. Comment corrigeriez-vous le cas de la racine double si vous connaissiez la multiplicité~?
\end{question}

\begin{question}[$\star\star\star\star$ \iqroles{QR, QT} \iqfirm{hedge fund} \iqlevel{confirmé}]\label{iq:iv1:calculus-and-optimisation:19}
Montrez que $\sum_{k \ge 1} 1/k^2 = \pi^2/6$, en calculant de deux manières l'intégrale double $\int_0^1\int_0^1 \frac{dx\,dy}{1 - xy}$. Indication~: développez en série géométrique pour une des méthodes, et effectuez une rotation des axes de 45~degrés pour l'autre.
\end{question}

\begin{question}[$\star\star\star\star$ \iqroles{QR, QT} \iqfirm{hedge fund} \iqlevel{confirmé}]\label{iq:iv1:calculus-and-optimisation:20}
Démontrez que $n! \sim C\sqrt{n}(n/e)^n$ pour une certaine constante $C$, en montrant qu'une suite convenablement choisie est monotone et bornée. Déterminez ensuite exactement la valeur $C = \sqrt{2\pi}$ à l'aide du résultat de Wallis, $\lim_{n \to \infty} I_{2n}/I_{2n + 1} = 1$, avec $I_n = \int_0^{\pi/2} \sin^n x\,dx$.
\end{question}
